\weeknum{7}
\topic[Fourier Transforms]{Fourier Transforms and Potential Fields}

\workshopstart

\vspace{-0.75em}
\noindent\textbf{Duration:} 3 hours \hfill \textbf{Setting:} Workshop room (laptops for demos)

\section*{Preparation}

\begin{description}
  \item[Pre-reading:] Course Notes, Ch. 17
  \item[Lecture videos:]
  \item[Notes:]
\end{description}

\section*{Purpose}

This workshop develops intuition for Fourier methods in potential-field geophysics
and applies them to gravity and magnetic interpretation.
The emphasis is on understanding how sampling, filtering, and phase
control what geological information can be recovered.

\section*{Learning Outcomes}

By the end of this workshop, students should be able to:
\begin{itemize}
  \item Relate spatial wavelength to source depth and geometry
  \item Recognize aliasing and understand its consequences
  \item Interpret gravity and magnetic spectra in geological terms
  \item Apply upward and downward continuation appropriately
  \item Explain the purpose and limitations of reduction to the pole (RTP)
\end{itemize}

\section*{Session Flow (\timelinelength\ min)}

\timeline[item]{Warm-up --- Two Bumps, One Line}{10}

Explore what happens when you have a more complex gravity anomaly.

\timeline[slot]{Mini-lecture --- Introduction to Fourier Transforms}{15}
\begin{itemize}
  \item Review of key ideas: wavelength, sampling, phase
  \item Why Fourier methods appear everywhere in potential-field geophysics
\end{itemize}

\timeline[slot]{Fold Belt Survey}{35}

Sampling decisions set a hard limit on geological resolution.  In this example, we'll explore the impact of aliasing.

\timeline[item]{Nevada Gravity: Filtering for Sediments}{50}
\begin{itemize}
  \item Complete Bouguer gravity of Nevada
  \item E--W and N--S power spectra
  \item Identification of Basin-and-Range wavelengths
  \item High-pass filtering and wavelength selection
  \item Parker--Oldenburg sediment thickness inversion
  \item Interpretation and limitations
  \item Effect of aliasing on inversion results
\end{itemize}

\timeline[slot]{Break}{5}

\noindent\textbf{Key emphasis:} Filtering encodes geological assumptions.

\timeline[item]{Tower Hill Magnetics: Righting the Anomaly}{40}
\begin{itemize}
  \item Magnetic anomalies as vector fields
  \item Cone/maar geometry and asymmetric anomalies
  \item Reduction to the pole as a phase operation
\end{itemize}

\noindent\textbf{Key emphasis:} RTP rotates phase; it does not filter noise.

\timeline[slot]{Wrap-up}{15}
\begin{itemize}
  \item When Fourier methods help---and when they mislead
  \item Connection to upcoming assignments and field interpretation
\end{itemize}

\timeline[slot]{Preview: Steady-State Heat Flow}{10}

Next week we move from potential fields to geothermics. We will develop a physical model for how heat moves through the crust under steady-state conditions, governed by Fourier's law of conduction. We will see how thermal conductivity contrasts refract isotherms in much the same way optical refraction bends light, and how heat production in the crust curves the geotherm. These ideas underpin both geothermal energy assessment and deep-Earth thermal modelling.

\section*{Key Ideas}

\begin{itemize}
  \item Fourier methods are powerful but unforgiving
  \item Sampling precedes interpretation
  \item Phase carries geometry
  \item Geological assumptions are embedded in every filter
\end{itemize}

\workshopend

\notepage
\notepage

\subimport{activities/}{activity_gravity_warmup}
\subimport{activities/}{activity_aliasing}
%\subimport{activities/}{activity_fft_demo1}
%\subimport{activities/}{activity_fft_demo2}
\subimport{activities/}{activity_nevada_fft_v2}
\subimport{activities/}{activity_tower_hill_rtp_v2}
